\section{EESM model and its three quadratic primitives}
\label{sec:model}

\subsection{Stationary model and conventions}

Let
\begin{equation}
 \gls{sym:i}=\begin{bmatrix}\gls{sym:id}&\gls{sym:iq}&\gls{sym:iexc}\end{bmatrix}^{\mathsf T},
 \qquad \gls{sym:dl}=\gls{sym:ld}-\gls{sym:lq},
 \qquad \gls{sym:kt}=\frac32\gls{sym:p}.
\end{equation}
For positive winding resistances and positive inductances, the stationary
linear model is
\begin{align}
 \gls{sym:pcu}&=\gls{sym:rexc}\gls{sym:iexc}^{2}
 +\frac32\gls{sym:rs}(\gls{sym:id}^{2}+\gls{sym:iq}^{2}),
 \label{eq:model-loss}\\
 \gls{sym:torque}&=\gls{sym:kt}\gls{sym:iq}
 (\gls{sym:lexc}\gls{sym:iexc}+\gls{sym:dl}\gls{sym:id}),
 \label{eq:model-torque}\\
 \gls{sym:ud}&=\gls{sym:rs}\gls{sym:id}
 -\gls{sym:omega}\gls{sym:lq}\gls{sym:iq},
 \label{eq:model-ud}\\
 \gls{sym:uq}&=\gls{sym:rs}\gls{sym:iq}
 +\gls{sym:omega}(\gls{sym:ld}\gls{sym:id}
 +\gls{sym:lexc}\gls{sym:iexc}).
 \label{eq:model-uq}
\end{align}
These equations follow the standard rotor-oriented synchronous-machine model
\cite{binder2017,schroeder2021}. Positive \(\gls{sym:iq}\) and positive
active flux give positive motoring torque; reversing either factor reverses
the torque.

\subsection{Matrix forms: entries before names}

The copper-loss matrix is
\begin{equation}
 \gls{sym:pcu}=\gls{sym:i}^{\mathsf T}\gls{sym:rmat}\gls{sym:i},\qquad
 \gls{sym:rmat}=\diag\!\left(\frac32\gls{sym:rs},
 \frac32\gls{sym:rs},\gls{sym:rexc}\right).
 \label{eq:R-matrix}
\end{equation}
Each diagonal coefficient is the watts paid for one squared ampere in that
winding coordinate. For \(\gls{sym:rs},\gls{sym:rexc}>0\), the signature is
\((+,+,+)\), the rank is three, and an equal-loss surface is an ellipsoid.
Its eigenvectors are the physical winding-current axes because this simplified
loss model contains no coupled loss term.

Torque is also quadratic, although it is visibly a product:
\begin{equation}
 \gls{sym:torque}=\gls{sym:i}^{\mathsf T}\gls{sym:hmat}\gls{sym:i},\qquad
 \gls{sym:hmat}=\frac{\gls{sym:kt}}2
 \begin{bmatrix}
 0&\gls{sym:dl}&0\\
 \gls{sym:dl}&0&\gls{sym:lexc}\\
 0&\gls{sym:lexc}&0
 \end{bmatrix}.
 \label{eq:H-matrix}
\end{equation}
The factor one half is required because each off-diagonal product appears
twice in \(\gls{sym:i}^{\mathsf T}\gls{sym:hmat}\gls{sym:i}\). Its eigenvalues
are
\begin{equation}
 0,\qquad \pm\frac{\gls{sym:kt}}2
 \sqrt{\gls{sym:dl}^{2}+\gls{sym:lexc}^{2}}.
 \label{eq:H-eigenvalues}
\end{equation}
Thus the generic signature is \((+,-,0)\): one direction raises torque, one
lowers it, and one is torque-neutral. A nonzero torque level is a hyperbolic
cylinder. At zero torque it collapses to the two planes
\(\gls{sym:iq}=0\) and
\(\gls{sym:dl}\gls{sym:id}+\gls{sym:lexc}\gls{sym:iexc}=0\).

The voltage map is linear:
\begin{equation}
 \gls{sym:u}=\gls{sym:amat}\gls{sym:i},\qquad
 \gls{sym:amat}=\begin{bmatrix}
 \gls{sym:rs}&-\gls{sym:omega}\gls{sym:lq}&0\\
 \gls{sym:omega}\gls{sym:ld}&\gls{sym:rs}&
 \gls{sym:omega}\gls{sym:lexc}
 \end{bmatrix}.
 \label{eq:A-matrix}
\end{equation}
Squaring its Euclidean output norm gives
\begin{equation}
 U^2=\gls{sym:i}^{\mathsf T}\gls{sym:qmat}\gls{sym:i},\qquad
 \gls{sym:qmat}=\gls{sym:amat}^{\mathsf T}\gls{sym:amat}
 =\begin{bmatrix}
 R_s^2+\omega_e^2L_d^2&R_s\omega_e\Delta L&\omega_e^2L_dL_{\mathrm{exc}}\\
 R_s\omega_e\Delta L&R_s^2+\omega_e^2L_q^2&R_s\omega_eL_{\mathrm{exc}}\\
 \omega_e^2L_dL_{\mathrm{exc}}&R_s\omega_eL_{\mathrm{exc}}&\omega_e^2L_{\mathrm{exc}}^2
 \end{bmatrix}.
 \label{eq:QU-matrix}
\end{equation}
Unadorned symbols in this displayed matrix denote the corresponding glossary
quantities and keep the coefficient audit legible. Because the matrix is a
Gram matrix, it is positive semidefinite. Because
\(\gls{sym:amat}\in\mathbb R^{2\times3}\), its rank is at most two. Under
ordinary nondegenerate conditions the signature is \((+,+,0)\): the voltage
boundary is an oblique elliptic cylinder, never a three-dimensional ellipsoid.
The null direction can be chosen as
\begin{equation}
 \vect n=\begin{bmatrix}
 -\omega_e^2L_qL_{\mathrm{exc}}\\
 -R_s\omega_eL_{\mathrm{exc}}\\
 R_s^2+\omega_e^2L_dL_q
 \end{bmatrix},\qquad \gls{sym:amat}\vect n=\vect0.
 \label{eq:voltage-null}
\end{equation}
Moving along this direction leaves both stator voltages unchanged. The zero
eigenvalue is therefore not an abstract degeneracy; it is the precise stator
current--field-current trade that is invisible to the inverter voltage in the
stationary model.

\begin{figure}[tb]
\centering
\begin{tikzpicture}
\begin{groupplot}[group style={group size=3 by 1,horizontal sep=4mm},
 width=.34\linewidth,height=4.8cm,publication 3d,
 ticks=none,xlabel={$i_d$},ylabel={$i_q$},zlabel={$i_{\mathrm{exc}}$}]
\nextgroupplot[title={loss}]
 \addplot3[surf,shader=interp,opacity=.45,draw=MidnightBlue!50,
 domain=0:360,y domain=-90:90,samples=19,samples y=11]
 ({1.05*cos(x)*cos(y)},{1.05*sin(x)*cos(y)},{.72*sin(y)});
\nextgroupplot[title={voltage}]
 \addplot3[mesh,opacity=.42,draw=BrickRed,
 domain=0:360,y domain=-1.2:1.2,samples=23,samples y=7]
 ({-.55*y+.72*cos(x)},{.50*sin(x)},{y});
 \addplot3[neutral axis,domain=-1.2:1.2,samples=2] ({-.55*x},{0},{x});
\nextgroupplot[title={torque}]
 \addplot3[surf,shader=interp,opacity=.42,draw=Purple!50,
 domain=.45:1.6,y domain=-1.1:1.1,samples=11,samples y=9]
 ({y},{.48/x},{(x-.35*y)/.9});
 \addplot3[surf,shader=interp,opacity=.42,draw=Purple!50,
 domain=-1.6:-.45,y domain=-1.1:1.1,samples=11,samples y=9]
 ({y},{.48/x},{(x-.35*y)/.9});
\end{groupplot}
\end{tikzpicture}
\caption{The three raw EESM objects are intrinsically different. Equal loss
is a closed ellipsoid; equal voltage is a rank-two elliptic cylinder with a
null axis; equal nonzero torque is a hyperbolic cylinder. Optimization is the
geometry of their first contacts.}
\label{fig:raw-geometry}
\end{figure}


\subsection{Why homogeneity changes the problem}

For every real scale \(a\),
\begin{equation}
 \gls{sym:pcu}(a\gls{sym:i})=a^2\gls{sym:pcu}(\gls{sym:i}),\quad
 \gls{sym:torque}(a\gls{sym:i})=a^2\gls{sym:torque}(\gls{sym:i}),\quad
 U^2(a\gls{sym:i})=a^2U^2(\gls{sym:i}).
 \label{eq:homogeneity}
\end{equation}
A direction selects the ratios \(M/P_{\mathrm{Cu}}\) and
\(U^2/P_{\mathrm{Cu}}\); radius selects the absolute level. Identifying
opposite nonzero vectors as the same direction is precisely the elementary
reason projective direction space is natural here. The sign of torque does
not disappear: \(\vect i\) and \(-\vect i\) give the same torque, while
changing only one torque-producing factor changes its sign.

This structure is lost in the ordinary two-current PSM equations because a
fixed magnet flux appears as an affine offset. It reappears when the PSM is
viewed as the plane \(\gls{sym:iexc}=\text{constant}\) inside the homogeneous
three-current parent problem.

\paragraph{Synthesis.}
Algebra gives three symmetric matrices. Geometry reads their signatures as an
ellipsoid, a hyperbolic cylinder, and an elliptic cylinder. Physics identifies
loss price, torque-producing products, and a voltage-invisible trade.
Computation gains a direction--radius separation. These claims require the
linear stationary model and positive winding resistances; later sections show
exactly what changes outside those assumptions.
